\section{Lexicon, Textual Variants, and Evidentiary Summary}
\label{app:register}

\subsection{Complete New Testament Occurrence Register of Key Burial and Textile Terms}

The table below lists every New Testament occurrence of the principal Greek textile and wrapping terms examined in the body, verified against the MorphGNT SBLGNT v6.12 lemma counts, the Robinson--Pierpont Byzantine Textform, and the Clementine Vulgate, alongside the Douay--Rheims (Challoner) English wording.

\Needspace{10\baselineskip}
\begin{fourstable}{Greek Lemma (Count)}{All NT Occurrences}{Clementine Vulgate}{Douay--Rheims Rendering}
othonion (5) & Luke 24:12; John 19:40; 20:5, 6, 7 & linteis (John 19:40); linteamina (John 20:5--7; Luke 24:12) & linen cloths (all 5 occurrences) \\
\midrule
othon\={e} (2) & Acts 10:11; 11:5 & linteum magnum & great linen sheet (10:11); great sheet (11:5) \\
\midrule
soudarion (4) & Luke 19:20; John 11:44; 20:7; Acts 19:12 & sudario (Luke 19:20; John 11:44); sudarium (John 20:7); sudaria (Acts 19:12) & napkin (Luke 19:20; John 11:44; 20:7); handkerchiefs (Acts 19:12) \\
\midrule
entyliss\={o} (3) & Matt 27:59; Luke 23:53; John 20:7 & involvit (Matt 27:59; Luke 23:53); involutum (John 20:7) & wrapped it up (Matt 27:59); wrapped him (Luke 23:53); wrapped up (John 20:7) \\
\midrule
sind\={o}n (6) & Matt 27:59; Mark 14:51, 52; 15:46 (2x); Luke 23:53 & sindone / sindonem & clean linen cloth (Matt); linen cloth (Mark 14); fine linen (Mark 15; Luke 23) \\
\midrule
keiria (1) & John 11:44 & institis & winding bands \\
\midrule
eneile\={o} (1) & Mark 15:46 & involvit & wrapped him up \\
\midrule
peride\={o} (1) & John 11:44 & erat ligata & was bound about \\
\midrule
parakypt\={o} (5) & Luke 24:12; John 20:5, 11; Jas 1:25; 1 Pet 1:12 & procumbens (Luke 24:12); inclinasset (John 20:5); inclinavit se et prospexit (John 20:11) & stooping down / stooped down \\
\midrule
ptyss\={o} / anaptyss\={o} (1 each) & Luke 4:20 (\latin{ptyxas}); Luke 4:17 (\latin{anaptyxas}) & cum plicuisset (4:20); ut revolvit (4:17) & when he had folded the book (4:20); as he unfolded the book (4:17) \\
\end{fourstable}

\Needspace{12\baselineskip}
\subsection{Textual Variants in John 20:3--10 and Luke 24:12}

\begin{fourstable}{Passage}{Byzantine Textform (RP2018)}{Critical / Alexandrian Witness}{Clementine Vulgate \& Notes}
John 20:6 & \latin{Erchetai oun Sim\={o}n Petros} & Nestle--Aland adds \latin{kai} before \latin{Sim\={o}n} (``Simon Peter also cometh'') & \latin{Venit ergo Simon Petrus} (no \latin{et} before \latin{Simon}) \\
\midrule
John 20:10 & \latin{pros heautous} (reflexive) & Nestle--Aland reads \latin{pros autous} (or \latin{hautous}) & \latin{ad semetipsos}; both forms carry the idiomatic sense ``to their home / lodging'' \\
\midrule
Luke 24:12 (whole verse) & Present in text & Westcott--Hort (1882, App. p. 71) bracket as a ``Western non-interpolation'' (omitted by D, Old Latin \latin{a b e r\textsuperscript{1}}, Eusebius's canons; present in \latin{c ff\textsuperscript{2}}, Old Syriac, Eusebius \emph{Ad Marinum}); SBLGNT prints in half-brackets & Present in full in the Clementine Vulgate; also attested by Papyrus 75, Vaticanus, and Sinaiticus \\
\midrule
Luke 24:12 (cloths) & \latin{ta othonia keimena mona} & SBLGNT omits \latin{keimena}: \latin{ta othonia mona} & \latin{linteamina sola posita} (retains ``laid'' / \latin{posita}) \\
\midrule
Luke 24:12 (departure) & \latin{pros heauton thaumaz\={o}n} & SBLGNT prints \latin{pros hauton thaumaz\={o}n} & \latin{abiit secum mirans} (attaches \latin{secum} to ``wondering''; RV 1895 and WEB-C attach \latin{pros hauton} to ``departed,'' yielding ``departed to his home'') \\
\end{fourstable}

\Needspace{12\baselineskip}
\subsection{Doctrinal, Exegetical, and Historical Grading Summary}

\begin{threestable}{Category}{Controlling Content}{Evidentiary / Authority Basis}
Dogma of Catholic Faith & Christ truly rose bodily on the third day by divine power, passing into immortal glory to die no more, and will come again in glory. & Creeds; Scripture (John 20; Luke 24; Rom 6:9; 1 Cor 15:20); Roman Catechism, Part I, Arts. III \& V. \\
\midrule
Direct Textual Facts (John 20:3--10; Luke 24:12) & Peter enters first and sees both \latin{othonia} and \latin{soudarion} rolled up (\latin{entetyligmenon}) apart (\latin{ch\={o}ris}) in one place; the beloved disciple enters second and ``saw and believed'' (singular), followed by ``they knew not the scripture'' (plural); Luke 24:12 shows Peter stooping, seeing the \latin{othonia} alone, and departing in wonder. & RP2018 Greek; SBLGNT Luke; MorphGNT SBLGNT lemma counts; Clementine Vulgate; Douay--Rheims; Thayer. \\
\midrule
Patristic and Scholastic Consensus (Anti-Theft Order) & The ordered separation of the myrrh-bound linen cloths and rolled-up head-cloth refutes hasty grave-robbery or confusion. & Chrysostom (\emph{Hom. in Io.} 85.4); Theophylact (PG 124); Aquinas (\emph{Super Io.} 20, lect. 1); Lapide (vol. 6). \\
\midrule
Legitimate Disagreement Among Fathers and Doctors & Whether ``and he saw, and believed'' (John 20:8) means Easter faith (Chrysostom, Cyril, Ambrose, Lapide for John alone) or belief in Mary's report of a taken body (Augustine, Gregory, Theophylact), with Aquinas presenting both side by side. & Chrysostom (\emph{Hom.} 85--87); Cyril (Book XII); Ambrose (\emph{In Luc.} X.179); Augustine (\emph{Tract.} 120.9, 121.1); Gregory (\emph{Hom.} 22.5); Theophylact (PG 124:291); Aquinas; Lapide. \\
\midrule
Spiritual and Liturgical Reception & The separated, rolled-up \latin{sudarium} figures God's transcendent eternity, unity, and Christ's guiltless Passion (Gregory); the Lazarus/Christ \latin{soudarion} contrast highlights immortal Resurrection; the Easter liturgy sings \latin{Angelicos testes, sudarium et vestes}. & Gregory (\emph{Hom.} 22.3--4); Augustine (\emph{Tract.} 49.24); Roman Catechism (Part I, Art. V); 1962 \emph{Missale Romanum}; 1981 \emph{Ordo Lectionum Missae} (no. 42). \\
\midrule
Bounded Historical Negative (Folded-Napkin Custom) & No examined Greek lexicon, rabbinic text (\emph{m. Berakhot} 8:3; \emph{b. Berakhot} 42a, 52b; \emph{m. Sanhedrin} 6:1), classical source (Smith; Suetonius), Hebraist commentary (Lightfoot), or patristic/scholastic author attests a first-century dining-napkin signal; earliest located circulation dates to 2006--2007 online texts. & Thayer; Mishnah; Babylonian Talmud; Smith; Suetonius; Lightfoot; patristic/scholastic commentators; \emph{Jerusalem Perspective} (2006/2013 capture); \emph{RecipesAndMore} (2007); \emph{TruthOrFiction} (2008/2015); \emph{GotQuestions}. \\
\midrule
Open Exegetical and Excluded Questions & Physical mechanics of how the cloths came to lie as found (unhurried removal vs. collapsed wrappings vs. Westcott's hybrid); Shroud of Turin and Sudarium of Oviedo forensics (excluded outside textual exegesis). & Henry Latham (\emph{The Risen Master}, pp. 34, 43); Brooke Foss Westcott (1908, vol. 2, p. 339); scope boundary. \\
\end{threestable}
